%% Introduction %%

Pour compléter notre propos, il est essentiel de comprendre que ces changements s'accompagnent d'une informatisation progressive des institutions patrimoniales et en conséquent du traitement documentaire. En effet, dans les années 70, la société s’informatise progressivement et l’INA ne fait pas exception. Elle se dote de ses premières bases de données et des questionnements documentaires qui les accompagne. 

\subsection{Histoire de la documentation}

%% Introduction %%

Pourtant, il est essentiel de ne pas négliger que ces questionnements sur le traitement et la politique documentaire n’ont pas attendu l’informatisation de la société pour être posée. Bien au contraire, ces nouvelles réflexions reposent sur des bases posées dès la naissance des livres imprimés. Pour comprendre les décisions prisent pour le traitement documentaire à la création de la cinémathèque puis à l'instauration du dépôt légal, il est important d'en comprendre les origines, les premières définitions et premières réflexions. Ils ont ouvert la voie à la mise en place d’outils et de techniques d’accès à l’information desquels les outils que nous utilisons aujourd'hui sont les héritiers. 

%% Naissance de la documentation %%

Ils naissent d’abord dans des cercles privés, avant de devenir collectifs\footcite[][15]{fayet-scribe2000}. Pour le constater, il nous suffit de nous plonger dans les origines de la BNF, quand à la révolution française Lefèvre d’Ormesson lance le projet de Bibliographie universelle de la France qui vise à constituer un inventaire des livres de la nation française, c’est à cette période que naît le terme « documenter », reprenant le sens ancien du terme « documenter » soit instruire et enseigner. A l’occasion de cette bibliographie nationale on y interroge également pour la première fois « l’ordre des connaissances »\footcite[][21-22]{fayet-scribe2000}, soit l’organisation des savoirs, comment réunir ces documents par les connaissances qu'ils renferment. 

%% Paul Otlet %%

En 1895 est crée à Bruxelles l’Institut international de bibliographie grâce à Paul Otlet et Henri Lafontaine. Ils sont considérés comme les créateurs de ce nouveau domaine qu’est la documentation\footcite[][45]{fayet-scribe2000} dont ils définissent le périmètre : 
\begin{displayquote}
	[…] des méthodes seront trouvées pour dépouiller les ouvrages avec suite et ensemble et permettre de retrouver instantanément, sans peine ni difficultés, les matériaux mêmes dont chaque publication fait apport à l’ensemble du savoir\footcite[Paul OTLET, « Les sciences bibliographiques et la documentation », Bulletin
	de l’IIB, 1903, pp.125-147. Retranscrit dans][48]{fayet-scribe2000}
\end{displayquote}
A cette occasion ils se séparent de la matérialité du manuscrit pour parler d’un « document », qui sert diffusent des informations, quelles que soient leur forme\footcite[][49]{fayet-scribe2000}. Cette définition trouve toujours un sens aujourd'hui. Nous utilisons au long de ce mémoire le terme général de \enquote{document} pour parler des images et sons conservés à l'Institut National de l'audiovisuel. Ces \enquote{documents} permettent de diffuser des informations, quelles qu'elles soient et nécessitent d'être organisés de façon à ce que tout un chacun puisse avoir accès aux connaissances qu'ils renferment. La documentation doit savoir s’adapter, peu importe la forme que prend le document et fournir toutes les informations permettant son identification et son partage\footcite[][55]{fayet-scribe2000}. Son projet se matérialisera en France à travers Charles Gariel et le général Hippollyte Sebert avec le Bureau bibliographique de Paris (BBP), branche française de l’Institut international de bibliographie\footcite[][62]{fayet-scribe2000}. Dans un objectif de normalisation, ils appliquent le classement décimal et naîtra en 1926 l’AFNOR, l’Association française de normalisation\footcite[][79]{fayet-scribe2000}. Paul Otlet et Henri Lafontaine poseront donc une base solide de l’institutionnalisation et l’internationalisation de la documentation des « documents ». Par cette internationalisation établissent l'importance de permettre le partage des description documentaire.

%% Suzanne Briet %%

A sa suite, Suzanne Briet, représentante de la Bibliothèque nationale au Bureau Bibliographique de Paris à partir de 1932. Elle constatera à la bibliothèque nationale que les outils documentaires mélanges nouveaux et anciens usages. En conséquent, elle créera un fichier d’information à l’attention des lecteurs afin de rassembler l’ensemble des catalogues, répertoires et bibliographies à leur disposition\footcite[][129]{fayet-scribe2000} et contribuera à l’ouverture de la salle des catalogues et des bibliographies à la bibliothèque nationale. Elle appuie ainsi l'importance de permettre aux chercheurs de comprendre ce que l'on pourrait appeler \enquote{l'architecture} du traitement documentaire auxquels ils sont confrontés. En ce sens, notre projet de cartographie du traitement documentaire de l'INA peut être rapproché de son projet. Tous deux ont pour volonté de donner toutes les clefs nécessaire aux chercheurs et plus largement aux utilisateurs des collection d'une institution patrimoniale pour réaliser la recherche la plus complète possible dans le catalogue de ces collection. En cela, le documentaliste est placé comme la clef du partage de ce savoir. Dans son manifeste « Qu’est-ce que la documentation ? », elle définit ainsi le métier de documentaliste : 
\begin{displayquote}
	1° est un spécialiste du fond, c’est-à-dire qu’il possède une spécialisation culturelle apparentée à celle de l’organisme qui l’emploie ; - 2° connaît les techniques de la forme des documents et de leur traitement (choix, conservation, sélection, reproduction); - 3° a le respect du document dans son intégrité physique et intellectuelle; - 4° est capable de procéder à une interprétation et à une sélection de valeur des documents dont il a la charge, en vue d’une distribution ou d’une synthèse documentaire.\footcite[][19]{briet}
\end{displayquote}
Définition qui plus de 70 ans plus tard entre toujours en résonance avec le métier de documentaliste actuel, malgré toutes les mutations que ce métier a connu. 

Nous distinguons déjà, dans ce même manifeste, les prémices de la transformation majeure que s’apprête à connaître le traitement documentaire : 
\begin{displayquote}
	A côté des fichiers et des catalogues qui présentent l’image schématisée des documents, par description abstraite de leur aspect formel, accompagnée ou non de photographie, on voit depuis peu de temps se constituer des fichiers parallèles obtenus par codification des éléments pouvant donner lieu à statistique ou à sélection. Le mot ici disparait, la lettre même est parfois absente, lorsque l’on a affaire aux machines à cartes perforées. La mécanographie statistique nous habitue à doubler les fichiers lisibles en clair par des fichiers où chacune des notations est la traduction conventionnelle des signes directement intelligibles. […] Le documentaliste sera de plus en plus tributaire d’un outillage dont la technicité augmente à une vitesse grand V. L’homo documentator doit se préparer à commander, toutes facultés en éveil, aux robots de demain\footcite[][28]{briet}.
\end{displayquote}

%% Rapprochement des disciplines %%

L’Histoire de la documentation et du traitement qui s’en accompagne est encore doté de plusieurs figures importantes et nous pouvons toujours y trouver des échos à la discipline que nous connaissons aujourd’hui. Tout au long de son évolution s’est toujours posé la problématique d’allier nouveaux et anciens outils. C’est une problématique à laquelle les institutions patrimoniales n’échapperont pas non plus. En 1993, Michel Melot, président du Conseil supérieur des bibliothèques évoque comme les disciplines propre aux bibliothèques, archives et documentation sont en train de se rapprocher grâce à la numérisation :
\begin{displayquote}
	La numérisation réduit tous les documents au même dénominateur commun, au même support physique. Nous entrons dans l'ère de l'unimédia. [...] 
	Bibliothèques et archives se rapprocheront. Quand les documents issus de la marche d'un service ou d'une entreprise seront disponibles sur les réseaux domestiques, la notion d'archives rejoindra celles de l'édition et lorsque l'éditeur sera le propriétaire de bases textuelles stockées sur serveur et accessibles à la demande, la notion d'édition aura rejoint celle des archives. Les accès en texte intégral auront, de même, fait se rejoindre les fonctions bibliothécaires (accès à l'unité bibliographique) et documentaires (accès au contenu des documents)\footcite[]{melot1993}.
\end{displayquote}
C'est en effet ce que nous pouvons constater en nous intéressant à la façon dont l'informatique s'est appropriée la documentation. 

\subsection{Prise en main de la documentation par l'informatique}

Nous pouvons reprendre notre développement aux années 70 et au début de l'informatisation de la société. Les bases de données commencent à être adoptées et comme l’avait introduit Suzanne Briet, ces « fichiers parallèles » s’imposent et les institutions s’adaptent. En 1968 émerge le format MARC (Machine-Readable Cataloging), permettant le stockage et l’échange de notices bibliographique par ordinateur qui sera rapidement adopté par de nombreuses bibliothèques, y compris la BNF. A l’INA, la première base de données est mise en place avant même sa création, par la vidéothèque de production. Nommée MIOR, elle garde toujours la structure de leur catalogue papier dans la continuité de ce mélange d’anciens et de nouveaux outils. Tout comme le système IMAGO en 1975, qui n'est pas une base de données à proprement parlé. Cette première version d’IMAGO est adaptée au thésaurus de l’INA et centralise toutes les collections qu’elle conserve\footcite[][42-44]{bassaïstéguy2025} mais elle reste toujours très attachée à la logique d’un catalogue de bibliothèque. Les notices étaient dans un premier temps saisies sur papier par les documentalistes puis saisies informatiquement par des prestataires. Deux fois par mois des index titres, générique et analytique étaient imprimées sur microfiche, ce qui se révéla néfaste pour certains champs du générique qui n'apparaissaient sur une seule émission d'une série pour gagner de la place\footcite[]{saintville1994}. La deuxième version d'IMAGO, en 1985 s’éloigne un peu plus des catalogues de bibliothèque grâce à différentes normes mises en place par la Fédération internationale des archives du film (FIAF), les \textit{Cataloging Rules for Film Archives}, pensées pour s’adapter à l’informatisation des catalogues de documents audiovisuels et permettre leur partage\footcite[][50]{bassaïstéguy2025} ainsi que par l'AFNOR avec les premières normes de l'image animée auxquelles l'INA a participé\footcite[]{saintville1994}. Nous pouvons souligner le parallèle avec l’importance que Paul Otlet donnait à l’internationalisation des normes de documentation. L’évolution des moyens informatiques y prend également une grande importance, les données sont transférées dans une base de données Mistral. Ce nouveau système permet cette fois la saisie en temps réel, la gestion informatisée des supports d'archives et la recherche documentaire\footcite{saintville1994}. Avec cette deuxième version, nous pouvons pour la première fois parler d'une base de données. 

%% Indexation %%

%En parallèle, les indexations évoluent en fonction des demandes professionnels. La demande de la première apparition d'une personnalité à la télévision revient souvent alors on ajoute l'indexation \enquote{première apparition}\footcite[]{saintville1994}, c'est une indexation adaptée à son public cible.  

%% Impact sur le traitement documentaire

Dans la chronologie des outils de consultation des documents conservés par l’INA, nous nous situons désormais à l’aube du dépôt légal et pourtant, l’arrivée de l’informatisation à l’INA génère déjà les premiers défauts dans les données de description. En passant d’IMAGO I à IMAGO II, certains champs ne correspondent plus. Quand les notices d’IMAGO I sont versées dans IMAGO II, les champs titres, sont automatiquement versés dans le champ titre collection, négligeant le titre de tranche horaire ou à un titre propre. Cela concerne potentiellement 230 000 documents réalisés avec IMAGO I. A la fin des versements, cela ne génère pas de problèmes lors de l’interrogation de la base de données car tous les champs titres sont interrogés\footcite[][53]{bassaïstéguy2025}. Mais aujourd’hui ce sont des défauts supplémentaires qui compromettent la recherche dans les bases de données, surtout dans la nouvelle base de données du dépôt légal. Les modalités de recherches dans cette base de données est tout à fait différente de celle de ce que l’on va désormais appeler les archives professionnelles. Les défauts effacés par cet outil de consultation apparaissent parfois bien plus clairement dans la base de données du dépôt légal. Et comme toutes les perturbations que nous avons déjà soulignées dans les données de descriptions dû à différentes politiques documentaires, la connaissance de ces évènements se perdent ou du moins ne se transmettent à chaque personne nécessitant d’effectuer une recherche dans une des bases de données de l’INA et l’accessibilité est largement compliquée.

%% Conclusion %% 

Malgré tout, cette informatisation, qui met à disposition de nouveaux outils pour la documentation de documents audiovisuels se place tout de même en héritier de cette longue histoire de la documentation. Le terme de document reste applicable aux nouveaux supports que conservent l'INA et la place des documentalistes reste essentielle à l'application des politiques documentaires. De même que ces politiques s'appuient toujours sur des normes favorisant le partage de ces documents. Ces outils commencent à particulièrement s'émanciper de leur héritage en s'adaptant aux besoins de leurs utilisateurs. La meilleure façon de mettre cela en avant est de nous intéresser à la création du nouvel outil de consultation et de production pour le dépôt légal. 